% What follows an \hfill in a paragraph or a heading is at the right of its
% line, as in the PDF. Also the places where an \hfill does nothing on the
% page.
\documentclass{article}
\usepackage{amsmath,amsthm,graphicx,subcaption,hyperref}
\begin{document}
Name \hfill Date

\hfill --- A. Author

First line \hfill right of the first\\
Second line \hfill right of the second\\
Third line, with nothing at its right

\textbf{Bold name \hfill bold date}

A paragraph long enough to fill its first line and wrap onto a second one,
which ends in a note at the right of that line. \hfill (a note)

a \hfill b \hfill c

\begin{enumerate}
    \item $f(x)=|x|$ \hfill (not differentiable at $0$)
    \item[\hfill x] A label of its own.
\end{enumerate}

\begin{description}
    \item[Term \hfill] Its description \hfill at the right
\end{description}

\begin{center}
    Centered left \hfill centered right
\end{center}

\begin{figure}[h]
    \centering
    \includegraphics[width=2cm]{test-picture.png}\hfill
    \includegraphics[width=2cm]{test-picture.png}
\end{figure}

\begin{tabular}{|l|p{4cm}|}
    \hline
    in a cell \hfill c & in a paragraph cell \hfill p \\
    \hline
\end{tabular}

\begin{figure}[h]
    \centering
    \includegraphics[width=2cm]{test-picture.png}
    \caption{A caption \hfill with a note at its right, long enough to take more than a line of the page, where the PDF sets it as a paragraph}
\end{figure}

\begin{figure}[h]
    \begin{subfigure}{0.45\textwidth}
        \centering
        \includegraphics[width=2cm]{test-picture.png}
        \caption{A subcaption \hfill with a note at its right, longer than a line}
    \end{subfigure}
\end{figure}

Boxes: \mbox{in \hfill an mbox} and \makebox[4cm]{in \hfill a makebox}.

Leaders: chapter \dotfill 5, and a rule \hrulefill\ here.

Before a break \hfill\break after the break.

\begin{proof}
Ends in text \hfill (by a lemma)
\end{proof}

\begin{center}
    \textbf{Centered above}\\
    Name \hfill Date
\end{center}

\begin{center}
    \href{https://example.com}{Linked left \hfill linked right}
\end{center}

\begin{flushright}
    Flush left \hfill flush right
\end{flushright}

Before a linebreak \hfill\linebreak after the linebreak, which starts the
next line at the left in the PDF.

Before a hint \hfill\linebreak[2] after the hint.

Before a penalty \hfill\penalty-10000 after the penalty.

\let\otherpenalty\penalty
Before another penalty \hfill\otherpenalty-10000 after it.

Short \hfill A remainder long enough to wrap onto a second line of the page,
which the PDF starts at the left of the line, as it does every line after the
first one of a paragraph.

\section{Heading \hfill at its right, or below it}
After the heading.

\begin{itemize}
    \item \href{https://example.com}{in a link \hfill linked}
\end{itemize}

\tableofcontents
\end{document}
